Die Arbeit an dem Projekt „GameUp“ im Rahmen meines Praxisprojekts war mit einer steilen Lernkurve verbunden. Da ich die alleinige Verantwortung für die technische Umsetzung des 
Gamification-Konzepts trug, war ich gefordert, mich in kurzer Zeit in neue Technologien einzuarbeiten und bestehendes Wissen auf ein professionelles Niveau zu heben.

\subsection{Vertiefung der Backend-Expertise}
Im Bereich der Backend-Entwicklung konnte ich meine Kenntnisse in \textit{PHP} und insbesondere im Umgang mit dem \textit{Doctrine Object-Relational-Mapper (ORM)} signifikant vertiefen. 
Die Herausforderung bestand hierbei nicht nur im Schreiben von Code, sondern im Entwurf einer Datenbankarchitektur, die komplex genug ist, um flexible Lerninhalte abzubilden, aber 
gleichzeitig performant genug für eine mobile Anwendung bleibt. Durch das Refactoring bestehender API-Endpunkte und die Implementierung der neuen Gamification-Logik habe ich ein tieferes 
Verständnis für die Funktionsweise moderner MVC-Frameworks wie Symfony entwickelt.

\subsection{Einarbeitung in moderne Frontend-Technologien}
Ein wesentlicher Teil meiner Arbeit fand im Frontend der Applikation statt. Dies stellte für mich in zweifacher Hinsicht Neuland dar:
\begin{itemize}
    \item \textit{Vue.js und TypeScript:} Vor Beginn des Projekts hatte ich noch nie mit dem Framework \textit{Vue.js} gearbeitet. Die Einarbeitung in die komponentenbasierte Architektur 
    sowie die reaktive Datenbindung erforderte zu Beginn eine intensive Auseinandersetzung mit der Dokumentation. Gleichzeitig eignete ich mir grundlegendes Wissen in \textit{TypeScript} 
    an. Der Wechsel von dynamischer zu statischer Typisierung im Frontend half mir dabei, die Codequalität zu erhöhen und Fehlerquellen bereits während der Entwicklung zu minimieren.
    \item \textit{Mobile App-Entwicklung:} Auch die Entwicklung von Applikationen für mobile Endgeräte war ein neuer Bereich für mich. Hier lernte ich die Besonderheiten im Umgang mit 
    mobilen Lebenszyklen (z. B. Ionic/Vue Lifecycle-Hooks) und die Berücksichtigung von Hardware-Ressourcen kennen.
\end{itemize}

\subsection{Entwicklung der Game-Engine als Lernerfolg}
Das technisch anspruchsvollste Teilprojekt war die Entwicklung der \textit{Game-Engine}. Hierbei habe ich weit über das normale Maß der Webentwicklung hinausgehende Konzepte gelernt:
\begin{itemize}
    \item Die Arbeit mit dem \textit{HTML5-Canvas-API} und die manuelle Steuerung von Rendering-Loops.
    \item Die Implementierung von mathematischen Berechnungen für die Positionierung und Interaktion von Entities.
    \item Die Abstraktion von Spiellogik in generische Szenentypen (\texttt{Slideshow\-Scene}, \texttt{Action\-Scene}), um eine automatisierte Inhaltsdarstellung zu ermöglichen.
\end{itemize}

\subsection{Fazit der persönlichen Entwicklung}
Rückblickend war das Praxisprojekt eine wertvolle Erfahrung, die meine Problemlösungskompetenz gestärkt hat. Besonders die Erkenntnis, dass ich in der Lage bin, eine komplexe technische 
Komponente wie eine Game-Engine von Grund auf eigenständig zu konzipieren und umzusetzen, hat mein Selbstvertrauen in meine Fähigkeiten als Softwareentwickler nachhaltig gefestigt. Die 
Zusammenarbeit im Team der snoopmedia GmbH gab mir zudem einen authentischen Einblick in professionelle Entwicklungsabläufe und die Bedeutung agiler Kommunikation.